\chapter{RESULTADOS E DISCUSSÃO}
\label{cap:resultados-discussao}

Os resultados do trabalho foram organizados em quatro entregas complementares: a revisão sistemática da literatura, a apreciação metodológica preliminar dos candidatos empíricos retidos, a consolidação dos fundamentos pedagógicos e a especificação conceitual do protótipo. Esta seção integra essas entregas e discute como elas responderam ao problema de pesquisa e aos objetivos definidos.

\section{Resultados da Revisão Sistemática}

A revisão identificou 11.904 registros em quatro fontes científicas. Após a remoção determinística de identidades DOI/URL e a triagem, 2.486 registros foram submetidos à priorização operacional baseada em metadados e resumos; essa etapa não corresponde à avaliação de elegibilidade de relatórios em texto completo. Dezoito registros foram retidos provisoriamente para adjudicação de escopo: 17 candidatos empíricos e 1 protocolo ou proposta contextual. A síntese empírica preliminar considerou apenas os candidatos empíricos e identificou aplicações voltadas à predição de desempenho e à estimativa de proficiência, incluindo modelos supervisionados como \textit{Random Forest}, SVM, árvores de decisão e redes neurais.

A heterogeneidade entre populações, instrumentos, variáveis e métricas impediu a comparação direta de todos os resultados. Acurácias reportadas em bases específicas foram interpretadas como evidências contextuais dos artigos, e não como garantias de desempenho em outros ambientes educacionais. Essa distinção evitou transferir resultados da literatura para o protótipo especificado sem validação empírica própria.

Os exemplos documentados na Seção~\ref{sec:tendencias-lacunas} incluem limitações de alocação e controle de confundimento, divergência entre duração prevista e período reportado, combinação de componentes de intervenção e sobreposição entre instrumento e variáveis preditoras. Eles ilustram problemas localizados em estudos consultados; a cobertura de textos primários e a ausência de uma matriz de codificação não permitem estimar a frequência dessas limitações na área.

\section{Resultados da Avaliação Metodológica}

A apreciação documental preliminar com o MMAT 2018 manteve as respostas aos cinco critérios separadas por estudo empírico. Esse formato tornou visíveis limitações relativas a representatividade, completude dos dados, controle de confundidores, integração de métodos e coerência entre desenho e análise, sem transformar os julgamentos em uma pontuação agregada.

A avaliação não produziu nota geral, média ou ranking de qualidade. Essa decisão preservou a interpretação de cada critério e impediu que uma pontuação agregada ocultasse fragilidades metodológicas relevantes. Como os julgamentos foram realizados por um único revisor e parte deles se baseou em resumos e metadados, seus resultados foram tratados como apreciação estruturada, auditável e preliminar, não como classificação definitiva dos artigos.

O fortalecimento da rastreabilidade também constituiu um resultado técnico. Os julgamentos foram registrados em arquivo canônico versionado e passaram a alimentar os artefatos CSV, HTML e LaTeX utilizados no trabalho. A correspondência entre avaliação e estudo considerou título normalizado, ano e primeiro autor, reduzindo o risco de associações incorretas ou ambíguas.

\section{Resultados da Fundamentação Pedagógica}

A fundamentação permitiu distinguir desempenho observado, proficiência estimada, competência e aprendizagem. Notas, acertos, tentativas e tempos de resposta representam evidências dependentes da tarefa e do contexto; a interpretação de conhecimentos e habilidades requer relacionar essas observações ao construto e ao instrumento de avaliação \cite{NRC2001}. Na especificação, o acompanhamento de mudanças ao longo do tempo foi tratado como condição para investigar aprendizagem, e não como resultado já demonstrado. Essa distinção também alcançou a probabilidade preditiva de um alvo definido: uma saída condicionada aos dados e ao modelo, distinta do indicador de desempenho observado e não equivalente, por si só, a proficiência, competência, aprendizagem ou eficácia pedagógica.

Os referenciais sobre conhecimento prévio, mediação, resolução de problemas, avaliação formativa e \textit{feedback} reforçaram que a utilidade de uma ferramenta computacional depende da forma como seus resultados são interpretados e incorporados à prática docente. A tecnologia foi posicionada como apoio à decisão, mantendo o professor responsável por considerar fatores pedagógicos, sociais e contextuais que não são observados diretamente pelo modelo.

Esses resultados fundamentaram requisitos de explicabilidade, revisão humana, associação curricular, proteção de dados e comunicação de incertezas. Assim, a dimensão pedagógica não foi adicionada apenas como justificativa; ela orientou diretamente a especificação técnica.

\section{Resultados da Especificação do Protótipo}

A integração entre revisão e fundamentação resultou em uma especificação conceitual composta por princípios de projeto, requisitos funcionais e não funcionais, critérios de seleção de fontes de dados, alternativas de modelagem, protocolo de avaliação e arquitetura de referência.

Entre os resultados consolidados destacaram-se a definição do protótipo como instrumento de apoio, sem substituição da interpretação docente, e a exigência de rastreabilidade entre dados, modelos, parâmetros e resultados. Também foram especificados modelos supervisionados simples e interpretáveis como alternativas iniciais de comparação, mantendo-se a separação entre avaliação técnica e alegações de eficácia pedagógica.

A especificação incorporou explicabilidade e análise de limitações ao fluxo de avaliação e organizou a arquitetura em componentes de ingestão, preparação, modelagem, avaliação, explicabilidade e apresentação. O uso de dados sintéticos foi delimitado a testes de fluxo, sem interpretação educacional.

A especificação constituiu um resultado final do trabalho porque transformou as evidências analisadas em decisões documentadas e verificáveis. Ao mesmo tempo, ela não foi apresentada como aplicação funcional nem recebeu métricas próprias de desempenho, preservando a correspondência entre o texto e os artefatos efetivamente produzidos.

\section{Atendimento aos Objetivos}

No escopo executado, o primeiro objetivo específico foi atendido pela categorização das técnicas e finalidades pedagógicas dos candidatos empíricos na síntese preliminar apresentada no Capítulo~\ref{cap:revisao}. O segundo foi atendido pela identificação de limitações e lacunas no material examinado, considerando a apreciação preliminar por critérios do MMAT e as restrições de acesso às fontes. O terceiro foi atendido pela formulação dos requisitos, critérios de dados e modelos e arquitetura de referência no Capítulo~\ref{cap:prototipo}. O quarto foi atendido pela proposição do protocolo de avaliação nesse mesmo capítulo; seu atendimento corresponde à definição dos procedimentos, não à execução de experimentos ou à demonstração de eficácia pedagógica.

O objetivo geral foi alcançado no escopo documental ao propor uma especificação técnica e pedagógica para apoiar a interpretação docente da aprendizagem matemática. A revisão e o \textit{pipeline} forneceram os meios e a rastreabilidade dessa elaboração. A especificação não foi validada empiricamente como solução educacional, e sua fundamentação permanece condicionada às limitações da síntese e da apreciação metodológica preliminares.

\section{Discussão Integrada}

A discussão da literatura e de suas limitações orientou a decisão de não selecionar uma técnica apenas pela frequência de uso ou pela maior acurácia reportada em um artigo. Na especificação proposta, a seleção deve considerar a qualidade e a representatividade dos dados, a finalidade pedagógica, a possibilidade de interpretação e as condições de aplicação.

A presença de modelos preditivos na síntese preliminar indicou tarefas de classificação e estimativa entre os candidatos analisados. Neste trabalho, essas tarefas não foram tratadas como mensuração completa de fenômenos educacionais. A especificação separou explicitamente desempenho, proficiência, competência e aprendizagem e previu participação docente na interpretação das saídas, sem atribuir aos autores dos estudos uma posição teórica uniforme.

A principal contribuição do trabalho esteve, portanto, na articulação entre revisão sistemática, avaliação metodológica, fundamentação pedagógica e engenharia de requisitos. Essa articulação ofereceu uma base estruturada para decisões futuras sobre soluções educacionais, mantendo a seleção de algoritmos subordinada ao contexto, à avaliação e aos limites da evidência.

%--------- FIM RESULTADOS E DISCUSSÃO ------------
